\section{Introduction}

The equivalent-layer technique has long being used in applied geophysics for 
for processing and interpreting potential-field data. 
It consists in approximating a discrete set of potential-field data 
by a linear combination of harmonic functions representing the potential field 
due to a discrete set of virtual sources (e.g., monopoles, dipoles, prisms).
These virtual sources are called \textit{equivalent sources} and are commonly disposed on a surface 
forming a discrete \textit{equivalent layer}, which is ``equivalent" 
in the sense that it satisfactorily reproduces the observed potential-field data.
The general approach involves estimating the constant physical property of the equivalent sources 
(e.g., density, mass, isotropic magnetic susceptibility, magnetization intensity, magnetic moment intensity) 
by numerically solving a constrained linear inverse problem.
The technique can be used to accomplish various data processing steps, such as interpolation, 
upward (or downward) continuation, gradient data processing, merging multiple data sets, 
estimating density models from gravity data, 
computing the amplitude of the anomalous magnetic field, 
transforming total-field anomalies into field components, 
as well as estimating the magnetization direction and computing the reduction to the pole
of magnetic data. 
Further details on the equivalent-layer technique can be found in the recent review by 
\citet{oliveirajr_etal2023}.

Most of methods based on the equivalent-layer technique 
are based on the widely accepted \textit{equivalent-layer principle}, which 
refers to the empirical generalizations by which most of the studies found in the geophysical literature justify 
the existence of equivalent layers allowing computing potential-field transformations.
In the companion paper ``Planar equivalent layers for potential fields: Part 1 -- Theoretical insights'',
we have proposed a comprehensive theoretical framework with the purpose of replacing the 
equivalent-layer principle with the \textit{planar equivalent-layer theorem}, 
which guarantees the existence of continuous and planar equivalent layers 
that reproduce the gravitational and magnetic induction fields due to 3D sources, 
even for the case in which the layer crosses them. 
From our theoretical developments, we have deduced different closed-form expressions for continuous and planar 
equivalent layers that are defined on the same horizontal surface to reproduce the a given potential-field data.

Here, we use the theoretical framework presented in the companion paper
``Planar equivalent layers for potential fields: Part 1 -- Theoretical insights" to explain 
well-established empirical approaches commonly found in the geophysical literature about the 
equivalent-layer technique, such as joint processing of gravity-gradient data 
\cite[e.g.,][]{barnes_lumley2011}, conversion of gravity disturbance data into the 
gravity-gradient tensor \cite[e.g.,][]{piauilino_etal2019_seg, oliveirajr_etal2023} and the use 
dipole layers to reduce magnetic data to the pole 
\citep[e.g.,][]{emilia1973, bhattacharyya_chan1977, silva1986, leao-silva_1989, oliveirajr_etal2013, 
li_etal2014, oliveirajr_etal2015, ribeiro_filho_etal2020, reis_etal2020, li-etal2021, takahashi-etal2022, piauilino_etal2025}. 
Results obtained from numerical simulations show that the estimated discrete equivalent layers 
commonly used in practical applications approximate the closed-form expressions we have presented in the 
companion paper. We also illustrate the application of an alternative equivalent-layer of dipoles to 
interpret airborne magnetic data over a complex alkaline province in central part of Brazil. Our results show that
the alternative approach can bring valuable information on the horizontal position of the true sources without 
imposing any information on their total-magnetization direction.
